 \documentclass{article}

%%%%% packages to import %%%%%%
\usepackage{graphicx} % Required for inserting images
\usepackage[useregional]{datetime2}
\usepackage[margin=1in]{geometry}
\usepackage{longtable}
\usepackage{subcaption}

%%%%% title information %%%%%%
\title{CS496 Software Project: \\ 
Web-Based HVAC Inventory App}
\author{Oscar Roat, Dylan Morales}
\date{\today}

\begin{document}
\maketitle % make title information appear here

\section{Client Information}
By sharing this client information and the rest of this document, you are stating that this client has provided this project as something they want (not something you created and asked if they wanted), and that they are interested in having you complete this project for your capstone.
% Complete the list items about your client
\begin{itemize}
    \item Client name: William Tikiob 
    \item Client title: Project Manager Engineer
    \item Client email address: william.tikiob@coolsys.com
    \item Client employer: CoolSys
    \item How you know the client: Introduced to by Dr. Isaacman
\end{itemize}

\section{Project Description}
% You must complete the following 4 subsections. The instructor will use this information to determine if your project might be feasible or not.

% Comment out the bracketed instructions as you finish subsections
\subsection{Overview}
%[Add a few paragraphs describing your project succinctly. What problem are you trying to solve, what is the purpose of your project? Why does your client want this project?]

The main priority of this software, and the problem we are trying to solve, is the tracking of materials for jobs. Tracking materials is divided among two roles. The project managers order the needed materials and keep track of the purchase orders (PO's) in files on a shared drive. At the job sites, the foremen receive the materials and sign off on the delivery. Currently, there is no way for the foremen to verify that all of the materials were actually shipped, which has lead to delays because of missing things. 

Our app would centralize the ordered materials in a place for everyone in the company to see. They would be organized by job, so comparing expected and delivered materials would be as simple as finding the appropriate job on the app. Project managers would be able to take photos of PO's, and - possibly using Optical Character Recognition (OCR) - the materials and their quantities would be inserted into a database. In addition, the people on site for these jobs should be able to take a photo of the packing slip, and the app would compare that list to the expected numbers in the database. 

Those on job sites will require an easy-to-use mobile app. It makes sense to make a system that will work both on a computer and on mobile devices. 

If the former is complete, there are secondary additions that the client requested. These are equipment serial capture and tool tracking. Tool tracking would allow all users to know what tools are currently allocated to certain people/jobs. Serial capture would allow the automation of inserting these tools into a DB. 

\subsection{Key Features}
%[At this point you should have a basic understanding of your client's needs. List out the key features of the software system the client wants you to build.]
\begin{itemize}
    \item Inserting material list into central DB (Project Manager)
    \item Organizing materials by job
    \item Comparing delivered materials to expected material list (Foremen)
    \item Automating insertion and comparison, possibly via OCR (client's recommendation)
    \item Custom role-based UI's (the simpler, the better)
    \item Mobile App design
    \item (stretch goal) Tracking tools by job and/or person 
    \item (stretch goal) Equipment serial capture
\end{itemize}

\subsection{Why this Project is Interesting}
%[Why did you decide this project was interesting enough to you to be a capstone project? What about this project is enticing? Why should anyone care?]
(Oscar) I like this project because of the high-level simplicity. The material lists are the primary data. All we want to do is provide a way for project managers to upload the lists and organize them, and for the foremen to review them. Obviously, it will be more difficult to implement than how I described it, but dealing with the logic of only one major data type is something that I am prepared to deal with. I also understand the perspective of the foremen and the confusion and anger that comes with incomplete orders on job sites. Being able to help their work run more smoothly is a bonus. 

\subsection{Areas of CS required}
[What subfields of computer science seem most likely to be relevant to your project? A capstone must involve multiple.]
\begin{enumerate}
    \item Software Engineering
    \item Web/Mobile App Development
    \item DB Management
    \item (possibly) Optical Character Recognition
\end{enumerate}

\subsection{Potential Concerns and Questions}
%[Is there any aspect of this project that makes you unsure if it will work, either due to your own interests/background, or that you aren't sure if it fits the requirements? Are there questions you have about this project that you want instructor feedback about?]
It is unclear whether or not we will be able to get to the secondary additions to the app. Going into the project, it seems that organizing and tracking tools should be similar to the material lists. Tool tracking could be moved from being a stretch goal to part of the main requirements. 
   
One concern we have is integrating OCR into the project. The company offered to pay for AWS Textract, which (according to them) has the necessary capabilities. This is interesting, but we are unsure of how this contributes to the difficulty of the project. While this is a concern, the focus at first should be on the presentation of the data. 

There are also a few questions we have, both for the client and in general:
\begin{enumerate}
    \item How would we get OCR to handle different PO formats, if necessary?
    \item How should we differentiate tools being used \textit{at a job site} versus \textit{by a specific person}?
    \item Should admin create accounts or should users create their own accounts?
    \item Who manages the tools (admin, pm, somebody else)?
    \item How do you handle delivery status?
\end{enumerate}

%\subsection{Summary of Efforts to Find a Project}
%(Not necessary for 482) [Briefly list out when/how you've discussed with this client, and if you've discussed with other clients who either didn't work out or didn't respond. If you considered a different project and it didn't work out, why didn't it work out?] 

%[Most CS495 projects end here. The sections below are for CS482 and CS496 software projects].

%\subsection{Comparison to Draft}
%[For CS496 only, focus on highlighting the major differences between the draft proposal in CS495 and this one here. If there are no major differences, you can remove this subsection.]

\section{Requirements}

\subsection{Non-Functional Requirements}
%[Non-functional requirements are just as important as functional requirements. Dont forget to specify them.]\\

Table~\ref{tab:nfr} presents the NFRs for this software project.

\begin{table}[h!]
\centering
\begin{tabular}{c l c p{9cm}}
\hline
\textbf{ID} & \textbf{NFR Title} & \textbf{Category} & \textbf{Description} \\
\hline
1 & Color Scheme & Visual / Design & No specifc colors required, but primary blue and red to match company colors \\ \hline
2 & Simple UI  & Visual / Design & The simpler the UI, the better \\ \hline
3 & Role-specific UI  & Portability & Different UI for different roles \\ \hline
4 & Mobile App  & Portability & Mobile version for foremen \\ \hline
2 & OCR Handling  & Reliability & OCR should account for different PO/Packing slip formats \\ \hline
\end{tabular}
\caption{Non-Functional requirements}\label{tab:nfr}
\end{table}


\subsection{Functional Requirements (User Stories)}
%[In CS482, all functional requirements are written as User Stories. In CS496, some projects may use a different template to write the requirements. The table below is an example of writing the Stories. Adapt accordingly to different templates or if you want to display more info.]\\

Table~\ref{tab:functional-req} presents the functional requirements for this software project.

\begin{longtable}[h]{c l c p{10cm}}
\hline
\textbf{ID} & \textbf{Story Title} & \textbf{Points} & \textbf{Description} \\
\hline
		
A1 & User Accounts & 8 & As a/an Admin, I want to create, update, and delete user accounts, so that each user has access to needed information \\ \hline
P1 & Material Lists & 3 & As a/an Project Manager, I want to upload material lists based on PO, so that users can see what is needed for projects \\ \hline
P2 & PO Photo & 3 & As a/an Project Manager, I want to upload a photo of a purchase order, so that I can avoid tidious input and errors \\ \hline
D1 & Projects & 5 & As a/an Project Manager, I want to create, update, and delete projects, so that I can keep projects up to date \\ \hline
D2 & Organize Material & 3 & As a/an Project Manager, I want to separate material lists by project, so that lists are organized \\ \hline
D3 & Assign Users &	2 & As a/an Project Manager, I want to assign users to specific projects, so that they are not distracted by other projects \\ \hline
D4 & Material List Status &	5 & As a/an Project Manager, I want to see the status of material deliveries across all projects, so that I can easily keep track of accuarcy and progress \\ \hline
D5 & Shipment Notification & 2 & As a/an Foreman/Logistics, I want a notification when materials are ordered, so that I know when to expect them at a project site \\ \hline
T1 & Tool Management & 8 & As a/an Foreman, I want to create, read, update, and delete company tools, so that I can manage them \\ \hline
D6 & Assign/Unassign Tools & 3 & As a/an Foreman, I want to assign tools to specific projects/people, so that users know what tools are in use and where \\ \hline
S1 & Record Shipment & 2 & As a/an Logistics, I want to input delivered materials , so that I can create a shipment and read a PO to compare for accuacy \\ \hline
S2 & Shipment Accuarcy & 5 & As a/an Logistics, I want to compare packing lists to material lists, so that I know if any materials are missing \\ \hline
S3 & Shipment Photo Upload & 3 & As a/an Logistics, I want to upload a photo of a packing slip, so that delivered materials can be automatically recorded \\ \hline
D7 & Tool Allocation & 2 & As a/an Foreman/Logistics, I want to see my tool when I am assigned to a project, so that other users know where a specific tool is \\ \hline
Q1 & Project Details on Forms &	2 & As a/an Quality Control, I want documentation to populate with project details, so that time is saved and error is minimized on projects \\ \hline
Q2 & Equipment Serial Capture &	5 & As a/an Quality Control, I want to take a photo of a serial nameplate and input its details onto QC form, so that time is saved and error is minimized on projects \\ \hline
M1 & Persistent Login and Logout & 3 & As a/an User, I want to login, remain logged in, and logout of my account, so that my information is secure and maintained \\ \hline
D8 & User Dashboard & 3 & As a/an User, I want to see my info (role, tools, projects) on my profile so that I have all of my personal info available \\ \hline
P3 & Project Popup & 3 & As a/an Project Manager, I want a popup of a project that I search, so that I can quickly see the information I want without other projects cluttering the screen \\ \hline
T2 & Tool Popup & 3 & As a/an Foreman/Logistics, I want a popup of the tool that I search, so that I can quickly see the information I want without other tools cluttering the screen \\ \hline
A2 & User Popup & 3 & As a/an Admin, I want a popup of the user that I search, so that I can quickly see the information I want without other users cluttering the screen \\ \hline


\multicolumn{2}{r}{\bf Total: } & 73 & \\ \hline

\caption{Functional requirements as User Stories.}\label{tab:functional-req}
\end{longtable}

\subsection{AI generated stories}
Given the project description and current user stories, Claude gave around 20 more stories, separated by category:

Authentication & Access
\begin{itemize}
    \item As a User, I want to reset my password, so that I can regain access to my account if I forget it.
    \item As an Admin, I want to assign roles to users, so that each user only has access to features relevant to their role.
    \item As an Admin, I want to view a list of all users and their roles, so that I can manage the team effectively.
\end{itemize}

Purchase Orders \& Materials
\begin{itemize}
    \item As a Project Manager, I want to edit or remove individual materials from a PO after it has been uploaded, so that I can correct OCR errors or changes in the order.
    \item As a Project Manager, I want to be notified when a delivery is flagged as incomplete or inaccurate, so that I can follow up with the supplier.
    \item As a Project Manager, I want to view the history of all POs associated with a project, so that I have a complete audit trail.
    \item As a Logistics user, I want to flag specific missing or incorrect materials on a delivery, so that discrepancies are clearly documented and actionable.
    \item As a Logistics user, I want to partially receive a shipment, so that I can record materials that arrived while others are still pending.
\end{itemize}

Notifications
\begin{itemize}
    \item As a Project Manager, I want to receive a notification when a packing slip has been reviewed and compared, so that I am aware of delivery outcomes in real time.
    \item As a user, I want to configure which notifications I receive, so that I am not overwhelmed by irrelevant alerts.
\end{itemize}
 
Projects
\begin{itemize}
    \item As a Project Manager, I want to mark a project as complete or inactive, so that it is separated from active projects without being deleted.
    \item As a User, I want to search or filter projects, so that I can quickly find the one I need.
\end{itemize}
 
Tools
\begin{itemize}
    \item As a Foreman, I want to mark a tool as lost, damaged, or retired, so that its status is accurately reflected for all users.
    \item As a Foreman, I want to view the full assignment history of a tool, so that I can track its usage over time.
    \item As a User, I want to search for a specific tool by name or ID, so that I can quickly find its current location and status.
\end{itemize}
 
Serial Capture / QC
\begin{itemize}
    \item As a Quality Control user, I want to view all serial numbers captured for a project, so that I have a complete equipment record.
    \item As a Quality Control user, I want to edit or correct a captured serial number, so that data entry errors can be fixed.
\end{itemize}
 
General / Mobile
\begin{itemize}
    \item As a User, I want the app to work offline or in low-connectivity environments, so that I can still log deliveries on remote job sites.
    \item As a User, I want the interface to be optimized for mobile, so that I can easily use it on-site without a computer.
\end{itemize}
 
Claude also gave feedback on what was missing in our stories: 
"The biggest gaps I noticed were around error correction (fixing bad OCR results), partial deliveries, discrepancy flagging with clear ownership, project lifecycle management, and tool status tracking beyond just assignment. Those are likely to come up early in development so worth prioritizing."

There are clearly some mistakes. For example, the two stories under General/Mobile are better classified as NFR's. Furthermore, the assigning user roles story is already built into our A1 story. For the most part, though, these stories are good additions. If I had to pick a few to add to the project, I would add setting projects as complete/inactive and letting users reset their passwords. 

\section{System Design}
% Reminder: Comment out the [bracketed instructions] as you finish subsections

\subsection{Architecture}
%[Which type of software architecture are you team following? Layered architecture, MVC, other? What are the main modules for your software?]
%[Main modules are not the same as Layers. If you adopted any form of layered architecture (MVC included), then your layers already group components based on responsibility. Therefore, for modules, think about semantically related components. For example, in a parking lot, I could have a User, Payment, Parking (Vehicles), Contact/Issues modules.]

Our team has chosen to utilize the MVC structured architecture to allow for obvious focus point on our development. Using MVC we can split up our work into 3 main categories: Model, View, and Controller. Using this style of structure will promote scalability, mendability, and maintenance as we move further down the road. It will allow for developers to handle different parts of the project, avoiding the conflict of waiting on another group member to finish before starting your work. When it comes to tools we will use for each part of MVC that will be explained in our Technology section. The module will be the main branch that handles our data that will contain important information for the game consistently and securely. The controller is 2 the link between the model and the view. It will take in requests data from the user and retrieve data from the model to display on the view. The view is the display of the structure and it allows for the user to see the information our software is displaying. Many benefits will be created from using this structure that will allow for efficiency, debugging, organization, modular development, team collaboration, and maintainable software.

\subsection{Diagrams}
%[CS482, on sprints/iterations 2-3, you need to create and update a diagram (check Iteration 2-3 assignment for which type of diagram). On CS496, since before sprint/iteration 1, you should have a class diagram and keep it up-to-date. In CS496, if your class diagram changes at each sprint, then create a Class Diagram subsection for the sprints, and show the changes; while keeping the one here the most up-to-date version. ]

\begin{figure}[b!]
    \centering
    \includegraphics[width=\linewidth]{Diagrams/Class1.jpeg}
    \caption{Current Class Diagram}
    \label{fig:currclassdiagram}
\end{figure}

See figure \ref{fig:currclassdiagram} for the current class diagram. I have a few notes for the current diagram:
\begin{itemize}
    \item Every job has one project manager and multiple logistics/foremen associated with it. I'm not sure if there's a way to differentiate them, since they are both 'Users,' so I wrote the appropriate role next to each arrow. 
    \item The arrow pointing from Job (should be 'Project') to Material references the Material List from the purchase order. The arrow pointing from Shipment Slip to Material refers to the list of actually delivered materials.
\end{itemize}

\subsection{Technology}
%[ Which technologies are you going to use to implement your project? This should include the chosen programming language, main frameworks/libraries, and database or data storage. Testing framework is essential here as well.]

For our project we have planned out some crucial technology that will allow us to produce a product to out client exactly as he wants. We plan to use java script for our main programming language, NodeJS for our server side run time environment, Express for the web application framework for building APIs, HTML/React/CSS for our front end framework to create a present and interactive view of the users of this app. We will focus on using react for our mobile app and HTML + CSS for our web app. For our testing we are going to use Jest to make proper tests to make sure our software work effectively. Our client asked for OCR to be used in our project to convert photos with text to database ready text. We plan to implement Tesseract.js or AWS Textract for this function of our web app. For our database we will user MongoDB to store all of our information in order to produce accurate real time information to our users. Also bcrypt for hashing our passwords for better security. 

\subsection{Coding Standards}
%[Are your team going to follow any coding standards? For example, using a naming convention for Database tables (like only singular lowercase names). Another example, only allowing code with unit tests and above 60\% coverage to be committed (good convention since testing is going to be evaluated). If you need inspiration to define your coding standards, the Extreme Programming approach has a set of coding, design, and test rules.]

NoSQL data structure

Users
    (
        (
        \_id: ObjectId,
        id\_user: 1,
        username: 'sample123',
        password: ('123'), \#hashed
        name: 'sample user',
        id\_tools: (()),
        role: 'ProjectManager'
        )
        ...
    )
Projects
    (
        (
        \_id: ObjectId,
        id\_manager: 1,
        id\_workers: (()),
        id\_materials: (()),
        location: 'sample'
        )
    )

Tables should have a capital first letter but following letters are lowercase

Backend

/models
/controllers
/routes
/middleware
/services
/tests
app.js

Consistent naming conventions

Modular architecture

Minimum 50\% test coverage before a commit

Make sure both partners can agree with new code before commit to git

Secure authentication practices like hashing password and page restrictions




\subsection{Data}
%[What is the main structure of your data? In SQL-like databases, this would be the planning of the main tables, their attributes, and interactions with other tables (basically an ER diagram). In NoSQL databases, this would be the main collections and general attributes of the JSON you will store in each collection.]

%[Tip to better find and write the data your system will need. Go back to your User Stories and for each one, think to yourself: which attributes/fields do I need to store for this to work?]

%[Tip 2. When a system has many different roles for people, those are usually done in a single User table/collection. Especially when they share many common attributes/fields.]
\begin{itemize}
    \item Material - name, quantity
    \item User - type (Admin, Project Manager, Logistics, Foreman), name, username/email, password, tools being used
    \item Tools - serial number, model, in use (bool)
    \item Project - location, PM, L/F on site, material list
    \item Equipment - serial number, model, ..., location
\end{itemize}

The '...' under the equipment data refers to attributes that I am not sure of yet. It could include different voltage levels, quality, or other things. Dealing with the equipment is still a stretch goal for now, so this is something we will worry about later. 

\subsection{UI Mocks}
%[Define and draw/sketch/code the main UIs your user will interact with in your software. Add your UI mocks here and a short caption about it. Do not forget about the main forms and CRUD UIs.]

\begin{figure}[h]
    \centering

    \begin{subfigure}[t]{0.3\textwidth}
        \centering
        \includegraphics[width=\linewidth]{UI Mocks/IMG_2820.jpeg}
        \caption{Jobs Page for Project Managers}
    \end{subfigure}
    \hfill
    \begin{subfigure}[t]{0.3\textwidth}
        \centering
        \includegraphics[width=\linewidth]{UI Mocks/IMG_2821.jpeg}
        \caption{User Profile Page}
    \end{subfigure}
    \hfill
    \begin{subfigure}[t]{0.3\textwidth}
        \centering
        \includegraphics[width=\linewidth]{UI Mocks/IMG_2822.jpeg}
        \caption{Tools Page for Foremen}
    \end{subfigure}
    
    \vspace{1em}

    \begin{subfigure}[t]{0.3\textwidth}
        \centering
        \includegraphics[width=\linewidth]{UI Mocks/IMG_2823.jpeg}
        \caption{Admin Page}
    \end{subfigure}
    \hfill
    \begin{subfigure}[t]{0.3\textwidth}
        \centering
        \includegraphics[width=\linewidth]{UI Mocks/IMG_2824.jpeg}
        \caption{Job Page for Logistics/Foremen}
    \end{subfigure}
    
    \vspace{1em}
    \caption{UI Mocks}
    \label{fig:UIMock}
\end{figure}

The color scheme will probably match the company colors, blue and red, though the client said that color was not a big deal.


\section{Iterations}

\subsection{Iteration Planning}
%% Reminder **CS482**: Update your iteration planning at each sprint
% Reminder: Comment out the bracketed instructions as you finish subsections
%[In CS496, you plan all iterations beforehand. In CS482, you update the planning here at each iteration. ]\\

Table~\ref{tab:it-planning} shows the iteration planning.

\begin{table}[h!]
\centering
\begin{tabular}{c l p{7cm} c c}
\hline
 & & & \multicolumn{2}{c}{\bf Points }  \\
\textbf{It.} & \textbf{Dates} & \textbf{Stories} & \textbf{Planned} & \textbf{Done} \\
\hline
1 & 01/29 - 02/10 & A1 User Accounts, P1 Material Lists, D1 Projects & 16 & 16\\ \hline
2 & 02/01 - 03/01 & D2 Organize Materials, D3 Assign Users, T1 Tool Management & 13 & ?? \\ \hline
3 & 03/01 - 04/01 & D6 Assign/Unassign Tools, S1 Record Shipment, S2 Shipment Accuracy & 10 & ?? \\ \hline
4 & 04/01 - 05/01 & D4 Material List Status, P2 PO Photo, S3 Shipment Photo & 10+ & ??\\ \hline
5 & 05/01 - 06/01 & D5 Shipment Notification, Q1 Project Details on Forms, Q2 Equipment Serial Capture & 09 & ?? \\ \hline
\multicolumn{3}{r}{\bf Total: } & 58+ & 65 \\ \hline
\end{tabular}
\caption{Iteration Planning for Incremental Deliveries}\label{tab:it-planning}
\end{table}

[Note that the original planning was limited because we did not have all of the details for the project at the time. The specific iterations will have more stories done.]

\subsection{Iteration/Sprint 1}
\subsubsection{Planning}
%[Which stories did you plan for this iteration/sprint. Add the total points for this plan. You can also explain the reason behind your planning, and what major feature(s) your team is focusing on delivering by completing these stories. You may use a table for a summary display of the planning, but elaborate in text more detail in your focus and feature plan.]
\begin{itemize}
    \item Oscar:
        \begin{enumerate}
            \item A1 User Accounts (8 Points)
        \end{enumerate}

        \item Dylan:
            \begin{enumerate}
                \item D1 Projects P1 Material list (8 points)
            \end{enumerate}

    \item Total Points Planned: 8+
\end{itemize}

For the first iteration, we focused on setting up the essential DAO's and their controllers, those being users and projects. Also, setting up basic UI should try to get done to keep frontend and backend on the same track. 


\subsubsection{Work Done}
%[Which stories did you complete in this iteration/sprint. Which ones did you partially complete? Who worked on which story? You may elaborate in paragraph(s) to add more detail about the work done.]

\begin{itemize}
    \item Oscar:
        \begin{enumerate}
            \item A1 User Accounts (8 Points)
        \end{enumerate}

        \item Dylan:
            \begin{enumerate}
                \item D1 Projects (8 points)
            \end{enumerate}

    \item Total Points Done: 8+
\end{itemize}

\textbf{Oscar:} I completed model and controller for users, which includes full CRUD functionality. Basic UI was also set up for the users, including a form to register and a list of users. While the UI is simple, there are still improvements that can be made on the style.

\textbf{Dylan:} I had planned to most of projects crud and and a portion of material list but decided to just complete the full crud for Projects. I have created two UIs which was one to handle create view and delete but also added a button to make updates to a desired project. I still see room for alot of improvement in how users are added to projects.

\subsubsection{Testing Coverage}
[Testing is very important. Show your coverage here. Is this coverage good enough? Explain why you think so. Is it not good enough? Explain a plan to increase the coverage. You may also elaborate on why some artifacts do not undergo much testing. If the testing changed from the last iteration, explain the reasons.]

\begin{figure}
    \centering
    \includegraphics[width=1\linewidth]{Iteration1Coverage.png}
    \caption{Iteration 1 test coverage report}
    \label{fig:placeholder}
\end{figure}

I think our coverage is done very well because it reaches close to perfect coverage in all fields. I think there is room for very little improvements just because we are almost near perfection. User controller test is missing some lines just because they arent very important for testing but in time could easily be covered in the next iteration. This is the first iteration of testing and cant be compared to anything in previous iterations.





\subsubsection{Retroespective \& Reflection}
%[What were the pitfalls, challenges, and issues you had in this iteration? How can you address them to improve the process in the next iteration? Did anything not go according to plan? Why so and how to avoid the same mistake? Write a personal reflection on what you learned in this iteration (even if a small technical thing like Database storage).]

\textbf{Oscar:} The only thing that I wasn't able to do this iteration was the testing for the functions in the script on the users admin page. I will have to make that up eventually. However, both the user DAO and controller files are above 80\% coverage, so I am still confident in the user functionality on the website. I also did some work on working with Bootstrap. While I used it in my software engineering project, I mostly copied from my teammates and did not really learn it. Learning Bootstrap will be something that I build on throughtout this project. 

\textbf{Dylan:} When it comes to how well the usability of the projects I think it could a little work on selecting users. I ran out of time and just created a simple method of just adding the user ID manually instead of having a drop down. I belive once the drop down is added I will not have to add much work the the project crud. I think I could have spent more time researching how we could possibly use a camera to input information from documents because I think this will be a very strenuous problem later on.  


\subsection{Iteration/Sprint 2}
\subsubsection{Planning}
%[Which stories did you plan for this iteration/sprint. Add the total points for this plan. You can also explain the reason behind your planning, and what major feature(s) your team is focusing on delivering by completing these stories. You may use a table for a summary display of the planning, but elaborate in text more detail in your focus and feature plan.]

\begin{itemize}
    \item D2 Organize Materials (3 points)
    \item P1 Material Lists (3 Points) may be 5
    \item D3 Assign Users (2 points)
    \item T1 Tool Management (8 points)
    \item M1 Persistent Login/Logout (2 points)
    \item Total points: 18 
\end{itemize}

Now that users and projects are set up, the next step is to link them by assigning users to projects. Other things that we are starting are setting up tools and adding login/logout (which also means hashing passwords and setting access permissions by role). While the total points at the time this is written (15) is low, M1 might end up being worth more points, and D2 might be combined with another story. Since we have created a full crud for materials[8] we need to also create the full crud for shipment as well. Once shipment is complete we will begin to implement the comparing feature that will compare a material list to a shipment then update its status.  

\subsubsection{Work Done}
%[Which stories did you complete in this iteration/sprint. Which ones did you partially complete? Who worked on which story? You may elaborate in paragraph(s) to add more detail about the work done.]

\begin{itemize}
    \item Oscar:
        \begin{enumerate}
            \item T1 Tool Management (8 points)
            \item M1 Persistent Login/Logout (2 points)
            \item part of (D6) Assign/Unassign Tools (1 point)
        \end{enumerate}

    \item Dylan:
        \begin{enumerate}
            \item D2 Organize materials (3 points)
            \item P1 Material Lists (3 or 5 points)
                
        \end{enumerate}

    \item Total Points Done: 18+
\end{itemize}

\textbf{Dylan:} For this sprint it was important that we created the create of crud for materials to make sure that projects could have a material list in order to make sure project managers can create and also foreman can view. During implementation i realized that installing the full crud would be a good idea. I needed to change a user story from organize the materials to just a full crud version. After creating the create function I went ahead an implemented all the other features including update and delete, completing the CRUD.
\textbf{Oscar:} I was mostly successful in this sprint. Login/Logout is totally done, which includes sessions and encrypted passwords. The Tool Management is very close to being done. Finally, I have added functionality to assign users to tools, though it might make more sense to have it the other way around. I also made small improvements and additions to the User stuff. 

In general, functionality for users and tools is up on the website. The only issue is an unknown error when trying to update anything. Even though an error shows up, the update still goes through. This is something that I will keep working on, but I am not too worried about right now because there is still more functionality to add. 

\subsubsection{Testing Coverage}
%[Testing is very important. Show your coverage here. Is this coverage good enough? Explain why you think so. Is it not good enough? Explain a plan to increase the coverage. You may also elaborate on why some artifacts do not undergo much testing. If the testing changed from the last iteration, explain the reasons.]

\begin{figure}
    \centering
    \includegraphics[width=1\linewidth]{CoverageSprint2.png}
    \caption{Iteration 2 Coverage}
    \label{fig:placeholder}
\end{figure}

The testing for our project we would say for our project is up to standard. Most of our functions achieve good branch and coverage function but could use a little work to get perfection. In order to improve our coverage, we should make sure that error handling is tested because that is where we lack certain function coverage. Also, for the Material controller, we noticed problems with conflicts with other testing since it contains User and Projects as attributes. So finding a fix to that should also ensure better coverage. Most of the testing has stayed the same and used similar patterns to previous testing, but adapted to handle new functions. 



\subsubsection{Retroespective \& Reflection}
%[What were the pitfalls, challenges, and issues you had in this iteration? How can you address them to improve the process in the next iteration? Did anything not go according to plan? Why so and how to avoid the same mistake? Write a personal reflection on what you learned in this iteration (even if a small technical thing like Database storage).]

\textbf{Dylan:} Some of the challanges I faced was dealing with matrial list's inclusion of users and projects. It was a bit of a learning curve dealing with interactions between classes but once I got it down it all worked out well. Another problem was testing because we mentioned previously, testing the material controller caused problems that I think were due to the tests being run all at once. I removed certain testing to fix it, but I believe a test runInBand might work. This process of iteration two went pretty well since it was very similar to the first but once I begin to integrate shipments and comparing for validation, I expect to run into some problems. I would say a valuable thing I learned this iteration is how to make multiple classes all work with each other. It was something that was a bit out of my mind until now, but I found that using an observer-type method allowed for a change in one place to affect all the others.

\textbf{Oscar:} I didn't have trouble with adding any of the functionality this iteration, but I am not set on the visuals yet. There is definitely more I can do to make it look better. However, the combo of not being a very artistic person and the need to prioritize makes it tough to try. I hope that I can add to the design little by little as the project goes on, but that will only happen if I put some effort into it. 
Also, so far I have been working on stuff that I am comfortable with, but soon we will start getting into unfamiliar territory, especially with the OCR, so there will be more trial and error with the process down the road. 
Finally, I met with the client last week and asked for the PO and shipment slips we will need to base our materials (and OCR) off of. It hasn't been necessary yet, so it is not a big deal, but it will be needed soon. 



\subsection{Iteration/Sprint 3}
\subsubsection{Planning}
%[Which stories did you plan for this iteration/sprint. Add the total points for this plan. You can also explain the reason behind your planning, and what major feature(s) your team is focusing on delivering by completing these stories. You may use a table for a summary display of the planning, but elaborate in text more detail in your focus and feature plan.]

\begin{itemize}
    \item (the rest of) D6 Assign/Unassign Tools (2 points left)
    \item S1 Record Shipment (2 points)
    \item S2 Shipment Accuracy (5 points)
    \item D8 User Dashboard (5 points)
    \item Total points: 14
\end{itemize}

\subsubsection{Work Done}
%[Which stories did you complete in this iteration/sprint. Which ones did you partially complete? Who worked on which story? You may elaborate in paragraph(s) to add more detail about the work done.]

\begin{itemize}
    \item Oscar:
        \begin{enumerate}
            \item (the rest of) D6 Assign/Unassign Tools (2 points)
            \item D8 User Dashboard (3 points, previously 5)
            \item (Spike) UI Change (2 points)
        \end{enumerate}

    \item Dylan:
        \begin{enumerate}
            \item S1 Record Shipment (8 points(Full Crud), previously 2(Just create function)
            \item S2 Shipment Accuarcy (2 points partial of the full 5 points)
                
        \end{enumerate}

    \item Total Points Done: 17
\end{itemize}

\subsubsection{Testing Coverage}
%[Testing is very important. Show your coverage here. Is this coverage good enough? Explain why you think so. Is it not good enough? Explain a plan to increase the coverage. You may also elaborate on why some artifacts do not undergo much testing. If the testing changed from the last iteration, explain the reasons.]

\textbf{Oscar:} I was able to finish the tests for error handling, so all of the code that I have written, aside from the js files accompanying the html pages, is at minimum 97\% coverage. 

\begin{figure}
    \centering
    \includegraphics[width=1\linewidth]{CoverageSprint3(1).png} %% NEED TO ADD
    \caption{Iteration 3 Coverage}
    \label{fig:placeholder}
\end{figure}

\subsubsection{Retroespective \& Reflection}
%[What were the pitfalls, challenges, and issues you had in this iteration? How can you address them to improve the process in the next iteration? Did anything not go according to plan? Why so and how to avoid the same mistake? Write a personal reflection on what you learned in this iteration (even if a small technical thing like Database storage).]

\textbf{Oscar:} I did not get as much done as I wanted to this sprint. I was at a tournament in South Carolina for most of spring break, so I was unable to do any work for about a week. However, there was time before my trip that I should have done more work. I finished the assigning/unassigning tools story, though more might get added to it depending on if we decide to change the data display on the website. I finished the User Dashboard, but it only ended up being worth 3 points right now. If anything else is added to it, extra points might be added as an additional story. I also put in a lot of time (hopefully a spike's worth) into fixing the design of the website, and I hope that most of the changes will be permanent. I also added error handling tests to the DAO's and controllers that I wrote, which is something that I had to catch up on. Overall, I thought I did okay this sprint given the time constraint, but I will need to pick up the pace for the remaining sprints. 

\textbf{Dylan:} This sprint was definitly the hardest sprint out of all of them because I found that I was tiring myself out on the front-end aspect of Shipments. After receiving feedback on our front end I went to find out how to make better ways of selecting instead of a dropdown. I tried multiple methods but finally found one that works best. Involves populating a drop down based off what the user types in a text box. After this I ran into a issue that involes our planning structure, which was how we store our materials. Once I finished shipment and tried to incorporate material checks, I realized that a Project should have Orders and the Orders should have Material lists. This way when creating a shipment you can just select a Project then that Project populates a Orders and then you can enter the received materials for proper Shipment accuracy checking. I ran into a lot of problems that stopped a lot of progress, so i decided to just focus on completing the full crud of Shipment instead of doing the shipment accuracy. I was still able to a bit of shipment accuracy adding to about 2 points worth. 

\subsection{Iteration/Sprint 4}
%[CS496 has 5 sprints. CS482 only has only 3 sprints (remove Iterations 4 and 5 from this doc if you are writing a doc for 482]

\subsubsection{Planning}
%[Which stories did you plan for this iteration/sprint. Add the total points for this plan. You can also explain the reason behind your planning, and what major feature(s) your team is focusing on delivering by completing these stories. You may use a table for a summary display of the planning, but elaborate in text more detail in your focus and feature plan.]

\begin{itemize}
    \item P3 Project Popup (3 points)
    \item T2 Tool Popup (3 Points) 
    \item A2 User Popup (3 points)
    \item S1 Record Shipment (2 points) [adjustments]
    \item S2 Shipment Accuracy (3 points)
    \item D2 Organize Materials (3 points)
    \item P1 Material Lists (5 points)
\end{itemize}

\subsubsection{Work Done}
%[Which stories did you complete in this iteration/sprint. Which ones did you partially complete? Who worked on which story? You may elaborate in paragraph(s) to add more detail about the work done.]

\begin{itemize}
    \item Oscar:
        \begin{enumerate}
            \item P3 Project Popup (3 points)
            \item T2 Tool Popup (3 points)
            \item A2 User Popup (3 points)
            \item Fixed error on updating
        \end{enumerate}

    \item Dylan:
        \begin{enumerate}
            \item S1 Record Shipment (2 points) [adjustments]
            \item S2 Shipment Accuracy (3 points)
            \item D2 Organize Materials (3 points)
            \item P1 Material Lists (5 points)
                
        \end{enumerate}

    \item Total Points Done: 22 
\end{itemize}


\subsubsection{Testing Coverage}

\begin{figure}
    \centering
    \includegraphics[width=1\linewidth]{CoverageSprint4.png}
    \caption{Iteration 4 Coverage}
    \label{fig:placeholder}
\end{figure}

Testing for this iteration went very well, and we were able to make some improvements that made our coverage a lot better. First, we had some failed tests, but after using --runInBand it fixed all our issues, and all tests passed. Tesing for Matrials was removed since it was not needed, and testing for Purchase Order was added. We still have a little imperfections but overall the testing came out very well and we are proud of our current standings.


\subsubsection{Retroespective \& Reflection}
%[What were the pitfalls, challenges, and issues you had in this iteration? How can you address them to improve the process in the next iteration? Did anything not go according to plan? Why so and how to avoid the same mistake? Write a personal reflection on what you learned in this iteration (even if a small technical thing like Database storage).]

\textbf{Oscar:} This sprint did not go as originally planned for me. Our original idea for displaying all of the data was in tables. However, we underestimated the amount of stuff we would need to display. There are going to be hundreds of projects and thousands of tools active at the same time, which would not work in a table format. Because of this, we had to come up with a new idea. Dylan added a search feature that searches by an attribute (email for users, serial number for tools, etc), and I added a popup with all of the data's information. Even though this was a new user story, there was no new functionality added. It also took me a little over a day to solve the bug that popped up when updating stuff. Going into the final iteration, most of the functionality that we wanted to add has been done. However, most of it is in the context of admin control. We will need to add less exhaustive features for the other roles. 

\textbf{Dylan:} This sprint was pretty much just making adjustments after the client revision. I had originally meant to just make all the UI of my cruds better fit for mass amounts of data and also look better. Although after being told the material list was not exactly what he wanted, I had to completely delete Materials Crud and restart with a Purchase Order Crud. This took alot of work because it was making a full new CRUD but also changing every other class that it affected. After I created the implementation of Purchase Order,s I was finally able to do the Shipment accuracy, so when logging a shipment, it would check for accuracy with the corresponding Purchase Order. This was a very frustrating Iteration just because I made a mistake and had to delete a lot of code from iteration 2. 

\subsection{Iteration/Sprint 5}
\subsubsection{Planning}
%[Which stories did you plan for this iteration/sprint. Add the total points for this plan. You can also explain the reason behind your planning, and what major feature(s) your team is focusing on delivering by completing these stories. You may use a table for a summary display of the planning, but elaborate in text more detail in your focus and feature plan.]

\begin{itemize}
    \item D5 Shipment Notification 
    \item Q1 Project Details on Forms 
    \item Q2 Equipment Serial Capture 
    \item D1 Projects (beautification and UI) + (view project addition)
    \item P3 Project Pop (Fixed homepage pop up)
\end{itemize}

\textbf{Oscar:} These were the stories we had originally planned, but going into this iteration, we will not have enough time to complete the OCR and serial capture stories. Because of that, I am going to focus on polishing the stories I have already completed, as well as finishing the separation of functionalities by roles.

\textbf{Dylan:} I had this planned out because I felt that moving on to new features that are not a necessity would not be the smartest move. Instead I went back and made sure that our website could be stronger at its current point. I would have like to add a email notification but I might just add it later. Instead I added a way to view projects that way its layed out better for users, and also made the Project/ProjectEdit JS alot cleaner and more functional.

\subsubsection{Work Done}
%[Which stories did you complete in this iteration/sprint. Which ones did you partially complete? Who worked on which story? You may elaborate in paragraph(s) to add more detail about the work done.]

\textbf{Oscar:} I did not complete any new stories. Adding OCR would have required us to refactor the DAO's related to purchase orders, shipments, and projects. Instead, I focused on abstracting functions and making sure that each role had the correct level of functionality. For example, the search bar and popup functions are used on many pages throughout the website. Now, they all call functions from one file ($search.js$ and $popup.js$, respectively), which has helped with readability and debugging. I would say that I have done around 8 points worth of work this sprint. While that is a little low, I am confident that all of the stories that I have previously completed are very polished. 

\textbf{Dylan:} Like Oscar I was not able to complete any new stories although I had done some research on how implement a email system from when shipments are logged. I was not able to actual put thin into the app but I might in future time because i think it is a nice touch. My main goal was to make a view project html allowing for anyone to look at a porject and its contents. I spent most of my time experimenting with new styles. I also fixed a problem with popups on the landing page. I also changed a good amount of Projects UI and Style to make it fit in better with the rest of the webiste and also added a search function. 

\subsubsection{Testing Coverage}
%[Testing is very important. Show your coverage here. Is this coverage good enough? Explain why you think so. Is it not good enough? Explain a plan to increase the coverage. You may also elaborate on why some artifacts do not undergo much testing. If the testing changed from the last iteration, explain the reasons.]

All of the tests pass when run in the individual test files, but there's an issue with running them concurrently that wasn't able to be solved. Other than that, the coverage is good across the board for all tests. 

\begin{figure}
    \centering
    \includegraphics[width=1\linewidth]{CoverageSprint5.png}
    \caption{Iteration 5 Coverage}
    \label{fig:placeholder}
\end{figure}



\subsubsection{Retroespective \& Reflection}
%[What were the pitfalls, challenges, and issues you had in this iteration? How can you address them to improve the process in the next iteration? Did anything not go according to plan? Why so and how to avoid the same mistake? Write a personal reflection on what you learned in this iteration (even if a small technical thing like Database storage).]

\textbf{Oscar:} I was unhappy that we weren't able to incorporate the OCR and serial capture, but I am not super surprised. Going into the project, Dylan and I agreed to prioritize making sure that the purchase orders, shipment slips, and tool stuff could be done manually before we added automatic input using OCR. I am very confident that we did that well. Especially after this current iteration, as most of it was spent polishing the website. We also agreed that the serial capture would be a stretch goal and communicated that to the client before the project began. While there was probably more polishing that I could have done, I was very busy outside of this class for the last few weeks. I am a bit worried that I did complete any new user stories, but the only ones remaining were things related to OCR. Overall, I am okay with my work done this iteration.

\textbf{Dylan:} Oscar and I did set goals when we first started and understood that we wouldnt be able to give the client all that he wanted. Although I think we did a very good job at giving him all the was important for our website to work for his needs. Our client wanted OCR and serial capture but I think we made the right decision on polishing instead of adding new bug prone features. Im upset that we couldnt give him everything that he wanted but Im happy with what we accomplished. I think that it looks nice and behaves well with user interaction. I dont think I got alot done this iteration but honestly adding more features on the last iteration would not be the best move.

\section{Final Remarks}

\subsection{Overall Progress}
%[Have you completed everything? If so, present evidence on how you brought value to your client, and the overall client satisfaction. Otherwise, estimate how much progress you done and how long it would take to finish this project. Be concrete about your progress, you know how many story points your software is, how many points you completed (this shows your progress). You also how many points your team delivers at each iteration, therefore you can estimate how many more iterations it would take to finish the leftover points (show the math).]

The only things we were not able to complete were related to the OCR/Photo Uploading and the Equipment Serial Capture. Those categories take up 4 stories (P2, S3, Q1, Q2) totaling 13 total points. From what Dylan has told me, some of the work he has done for purchase orders and shipment slips would have to be refactored to work with OCR. In total, it would probably be anywhere from 17-20 points worth of work to integrate the remaining stories, which would take one more sprint if we were soley focused on that. Excluding those stories from our completed work, we finished 60 points worth of stories. While that is on the lower end in terms of points, I am confident that the completed stories are well done and will bring value to the client. 

\subsection{Project Reflection}
%[Your personal reflection on the project. What lessons did you learned. What would you have done differently? How can you do better work in future projects? You may write this as a team or per person (or both --- if all your iterations were team reflections, then it would be better to write individual reflections here)]

\textbf{Oscar: } Overall, my process for this project was smoother than that of my CS482 project. The primary reason was because I did not have to learn to use any new languages or tools; I was already familiar with the database (Mongo), server connection (Node, Express), and testing (Jest) tools> I was also much more familiar with working in Java script and HTML compared to before. The other reason for this process being smoother was my time management. In CS482, I would consistently wait until the last 3-4 days of each iteration to get most of my work done, which led to increased stress. Going into this project, I knew that things wouldn't perfectly according to plan and that stories would need to be adjusted based on client needs. However, as long as I remained active, working on the project almost every day, I could adjust accordingly. 

One thing that I could have done better was get better info about the requirements from my client. The main function of this website is to organize the active projects and company tools. Part of that organization includes being able to easily view any given project/tool. However, I grossly underestimated the amount of stuff the company would be managing. Our original idea of organizing data in tables had to be scrapped after it was implemented over the course of the first two sprints. That change slowed progress in iterations 3 and 4 and probably led to our inability to complete the OCR stories. If I had known the quantity of data from the start, a lot of time would have been saved. 

The biggest skill that I improved on was UI. My first design for the webpages looked very blocky and unappealing. I realized that I was not happy with it; I wanted the website to look as good as it ran. That is something that I had never prioritized before. I had to actively think about design more than I ever had, and I think that I ended up doing a good job. The current design now is not great, but it is a major improvement over the first design. 

Finally, I believe that I did a good job at maintaining the quality of our software. I always made sure to add appropriate testing at the same time that new functionality was added, with a few exceptions, so I wouldn't fall behind. Furthermore, I did much more refactoring than I did in CS482. Specifically, I found ways to abstract functions that were being used in multiple locations on the website. This made it easier for both me to debug and (hopefully) for the client to maintain in the future. 

\textbf{Dylan:}
I think this project went pretty well and I would say that we covered all the important parts of the project. Through out the project I always wish that we had planned better stories, beacuse will working I would run into to things that were needed but never mentioned. Which would lead to refactoring of earlier iterations. Also make sure you start off early with scalable methods like search instead of drop down menus or like popups instead of tables or cards. I found how important it is to constantly check with your client before you go too far in the wrong direction because that lead to some set back in out later iterations.

When it comes to doing things better I would say it would mainly be my forward thinking. Many times through out the project I ran into issues with my previous software. One main problem was making a materials crud when I should have just made a purchase order crud with matrials inside of it. We planned out stories to cover everything we need but didnt really think deep into the functionally of all the cross references. I also think that I should have starting working from the first day of each iteration instead of wating until the last week or couple days. This would make get very stressed and produce code that I wasnt too happy with.

For later work I will make sure to take everything in step but also carefully calculate how my step will work with my future ones. I also want to make sure I set up building block at the start of my future projects. Oscar seemed to do this very well but I think i lacked in that aspect. When I would make something that could be reused I would sometimes forget and just hard code it in. I think that for my next project I will just overall have to make sure I work on setting up a better foundation that way I can smoothly go through iteration with a strong plan and never second guess.

Overall, for this project I think I had learned a lot of new things and have enhanced how well i understand project planning and working with a team. I would say that when it comes to confidence with our app I would say im alot more confident with this project compared to CS482 because I had also tried to make sure everything functioned properly because it will real world use. This also has direct correlation with our testing since it shows alot better coverage then out CS482 testing. The actual UI of the project is not that sophisticated but I think it works exactly how our client would want. I would have liked it to look better but Im more focused on how well it work then compared to beauty. I do belevie that after making this project that any project in the future that I do will get done alot quicker and better.




\section*{Appendix}
[Appendix section if needed]


\end{document}
